\section{Harmonic coordinates and the quadratic estimate}\label{sec:harmonic}

Let $E=\bigoplus_i F_i\otimes W_i$ be a polystable bundle on
$S$, with pairwise nonisomorphic stable factors and a direct-sum
Hermitian--Einstein metric $h_E$. Put
\[
 V=H^1(S,\End E),\qquad
 G=\prod_i\GL(W_i),\qquad K=\prod_i\U(W_i).
\]
We use unnormalized trace inner products and volume form
$\omega^2/2$. Write $d=\dbar_E$, and let $H_j,G_j$ denote
harmonic projection and the Green operator in degree $(0,j)$.
The flat metric on $V$ is the real metric of this harmonic
Hermitian form. Its quaternionic structure is induced by the
parallel holomorphic symplectic form of $S$, normalized to make
that structure antiunitary. All quotient metrics use these
same normalizations.

\subsection{Exact quadraticity in the harmonic coordinates}
The canonical harmonic recursion is
\begin{equation}\label{eq:recursion}
 a(x)=x-d^*G_2(a(x)\wedge a(x)),\qquad
 \kappa_E(x)=H_2(a(x)\wedge a(x)).
\end{equation}
It converges on a neighbourhood of $0\in V$ and gives a
holomorphic semiuniversal deformation over
$Z=\kappa_E^{-1}(0)$; its underlying reduced space suffices
for all metric statements. In particular,
\begin{equation}\label{eq:coclosed}
 H_1a=x,\qquad d^*a=0,\qquad a=x+O(|x|^2).
\end{equation}
The maps are locally $G$-equivariant. These assertions, including
integrability precisely on $Z$, are the harmonic construction
of \cite[Lemma~2.1]{wuebben15}. Convergence follows from the
contraction for $d^*G_2(a\wedge a)$ in a sufficiently high
Sobolev norm. Parallel endomorphisms commute with the Hodge
operators, which gives the equivariance even for nonunitary
elements of $G$ whenever both sides are defined.

\begin{lemma}\label{lem:quadraticity}
On the full convergence neighbourhood of \eqref{eq:recursion},
\begin{equation}\label{eq:exact-quadraticity}
 \kappa_E(x)=H_2(x\wedge x).
\end{equation}
\end{lemma}

\begin{proof}
Choose a nonzero parallel holomorphic two-form $\sigma$ on $S$.
Every endomorphism-valued $(0,2)$-form is $\phi\bar\sigma$.
The formula for $d^*$, together with parallelness of $\sigma$,
identifies $d^*(\phi\bar\sigma)$ with a nonzero constant
multiple of contraction of $\partial_E\phi$ by $\bar\sigma$.
Thus $\phi\bar\sigma$ is harmonic if and only if
$\partial_E\phi=0$. Equivalently, $\phi^*$ is holomorphic.
The induced connection on $\End E$ has zero Hermitian--Einstein
constant, so its holomorphic sections are parallel. We obtain
\begin{equation}\label{eq:harmonic-two}
 \mathcal H^{0,2}(\End E)
   =\{\xi\bar\sigma:\nabla\xi=0\}.
\end{equation}
Consequently, for $\theta=\theta_0\bar\sigma$,
\begin{equation}\label{eq:harmonic-test}
 H_2\theta=0\quad\Longleftrightarrow\quad
 \int_S\tr(\theta_0\xi^*)\,dV=0
 \quad\hbox{for every parallel }\xi.
\end{equation}

We claim that for $d^*u=0$ and arbitrary
$\gamma=\phi\bar\sigma$,
\begin{equation}\label{eq:coexact-vanishing}
 H_2(u\wedge d^*\gamma)=0
      =H_2(d^*\gamma\wedge u).
\end{equation}
Use unitary normal frames at a point. Up to fixed nonzero
scalars, the two coefficients relative to $\bar\sigma$ are
\[
 -\sum_j u_{\bar j}\nabla_j\phi,
 \qquad \sum_j(\nabla_j\phi)u_{\bar j}.
\]
For a parallel $\xi$, covariant integration by parts gives
\begin{align}
 \int_S\sum_j\tr(u_{\bar j}\nabla_j\phi\,\xi^*)\,dV
 &=-\int_S\tr\left(\left(\sum_j\nabla_j u_{\bar j}\right)
                   \phi\xi^*\right)dV=0,\label{eq:ordered-one}\\
 \int_S\sum_j\tr((\nabla_j\phi)u_{\bar j}\xi^*)\,dV
 &=-\int_S\tr\left(\phi\left(\sum_j\nabla_j u_{\bar j}\right)
                   \xi^*\right)dV=0.\label{eq:ordered-two}
\end{align}
The contracted derivative is a constant multiple of $-d^*u$.
The connection preserves the metric and volume, and
$\nabla\xi^*=0$, so these are global divergence identities.
The order of the endomorphisms is retained in both formulas;
neither commuting factors nor flatness of the Chern connection
is used. Equations \eqref{eq:harmonic-test}--\eqref{eq:ordered-two}
prove \eqref{eq:coexact-vanishing}.

Now write $a=x+\beta$, where
$\beta=-d^*G_2(a\wedge a)$. Both $x$ and $\beta$ are
$d^*$-closed. Formula \eqref{eq:coexact-vanishing}, first with
$u=x$ and then with $u=\beta$, annihilates under $H_2$ all
three terms $x\wedge\beta$, $\beta\wedge x$ and
$\beta\wedge\beta$. Expanding $H_2(a\wedge a)$ proves
\eqref{eq:exact-quadraticity}.
\end{proof}

This calculation specifies the coordinates in the quadraticity
theorems for hyperholomorphic bundles. It uses the central
Ricci-flat metric on the surface. The coordinate specification
is essential for the metric estimate below: a general
holomorphic substitution $x\mapsto x+O(|x|^2)$ has derivative
$I+O(|x|)$ and need only preserve a relative linear error.

Choose unitary harmonic bases for the cross-extension spaces
and symplectically paired bases for the loops. Serre duality
identifies $H_2(x\wedge x)$ with the complex moment map of this
doubled quiver. More precisely, with $r_i=\rk F_i$ and a
normalized harmonic generator $\eta_S$ of $H^2(S,\OO_S)$,
\begin{equation}\label{eq:rank-dual}
 H_2(x\wedge x)
   =\bigoplus_i r_i^{-1}I_{F_i}\otimes
                      \mu_{\CC,i}(x)\,\eta_S,
\end{equation}
up to the single common nonzero symplectic normalization;
see the rank-dual formula in \cite[Section~3]{wuebben14}.
In particular the zero sets agree exactly. Since
$H^1(S,\OO_S)=0$, fixing the determinant removes no harmonic
direction. Lemma~\ref{lem:quadraticity} therefore gives
\begin{equation}\label{eq:Z-cone}
 Z=\mu_{\CC}^{-1}(0),\qquad
 (Z\cap\mu_{\RR}^{-1}(0))/K=\Cone_E.
\end{equation}

\subsection{The quantitative local comparison}
We record precisely the metric input from
\cite[Theorem~1.1(i),(iii)]{wuebben15}, at constant polarization.
For a fixed locally free polystable background on a compact
K\"ahler manifold, in the coordinates \eqref{eq:recursion} and
on the fixed-determinant subspace, that theorem gives a
homeomorphism
\[
 f_0:(Z\cap\mu_{\RR}^{-1}(0))/K\longrightarrow
       \text{a neighbourhood of }[E]
\]
onto the nearby Hermitian--Einstein moduli. It covers every
sufficiently nearby solution, preserves the local complex-orbit
description, and is smooth on smooth pieces of fixed orbit type.
For genuine tangent vectors to such pieces,
\begin{equation}\label{eq:imported-estimate}
 (1-C|x|^2)g_{\rm fin}
 \le f_0^*g^{L^2}\le(1+C|x|^2)g_{\rm fin}.
\end{equation}
Here, if $v$ is a balanced tangent lift and $C_x\xi=[\xi,x]$,
\[
 g_{{\rm fin},x}([v],[v])=\inf_\xi\|v-C_x\xi\|^2
\]
is the flat reduction metric; the scalar kernel can be omitted
from the gauge algebra. This is a statement on actual stratum
tangents, not on all Zariski tangents at a singular point.
The constant is independent of stabilizer type, of the positive
singular values of $C_x$ or $D\kappa_E(x)$, and of ratios of
positive orbit eigenvalues. It may depend on the fixed
background, ranks and multiplicities, harmonic dimensions,
positive elliptic gaps, Sobolev constants and the chosen
regularity order. These are exactly the dependencies needed
in \eqref{eq:main-comparison}.

All hypotheses hold here by Proposition~\ref{prop:no-bubbling}.
The Hermitian--Einstein correspondence identifies the analytic
moduli with the slope-polystable bundle moduli, and the latter
is the Gieseker moduli germ by that proposition. By
\eqref{eq:Z-cone}, $g_{\rm fin}=g_{\Cone_E}$. A zero-level
balanced representative satisfies
\begin{equation}\label{eq:cone-radius}
 r([x])=\|x\|.
\end{equation}
Indeed, the radial segment is horizontal and has length
$\|x\|$, while the norm is $1$-Lipschitz on every stratum.
The full cone distance has the same property
\cite[Section~5]{wuebben14}. Thus \eqref{eq:imported-estimate}
is precisely the asserted quadratic tensor estimate.

It remains to check that the identification is holomorphic in
these same coordinates. The family $\dbar_E+a(x)$ is
holomorphic, locally $G$-equivariant and semiuniversal. The
analytic GIT construction of Buchdahl--Schumacher
\cite[Theorems~5--9 and Remark~4]{buchdahl21} identifies its
reduced analytic quotient $Z/\!/G$ with the reduced moduli
germ. Denote its holomorphic classifying isomorphism by $\chi$.
The Kempf--Ness identification supplies the homeomorphism
between the balanced quotient in \eqref{eq:Z-cone} and
$Z/\!/G$.

The map $f_0$ in \eqref{eq:imported-estimate} corrects the
same operator $\dbar_E+a(x)$ within its complex gauge orbit.
Its finite balancing step changes $x$ within its $G$-orbit.
Consequently $f_0$ and $\chi$ assign the same bundle class to
every balanced representative. Local coverage and uniqueness
identify them. The holomorphic structure on the cone quotient
therefore makes $\Psi_E=f_0$ a biholomorphism of germs, without
any further change of coordinates. This proves part~(ii) of
Theorem~\ref{thm:main}.

\subsection{Finite-dimensional slices used below}\label{subsec:slice-input}
We also need two finite-dimensional conclusions of
\cite[Sections~4--6]{wuebben14}. At a fixed point of a flat
zero-level doubled-quiver quotient, its full quaternionic
normal representation $W$ is again a doubled Ext-quiver
representation. If $K'$ is the stabilizer and $U=W^{K'}$,
the normal quotient splits as the exact metric product
\begin{equation}\label{eq:normal-product}
 W\hkr_0K'=U\times\bigl(U^\perp\hkr_0K'\bigr).
\end{equation}
The stratum through the centre is $U\times\{0\}$. There is a
slice homeomorphism, smooth on strata, with relative metric
error $O(\|w\|)$ and derivative at the central stratum equal
to the isometric normal identification. Its constant may
depend on that centre, but not on positive eigenvalues of the
retained orbit operator.

The quotient is proper and geodesic, with its quotient
topology. If its complex zero fibre has one module-simple
point, then the simple locus is dense and every full normal
Ext-quiver has a module-simple point as well. These are the
metric-slice theorem, the properness theorem and the
normal-simple lemma of \cite{wuebben14}. The explicit simple
points in Section~\ref{subsec:quivers} verify this hypothesis
for our two singular cones. These statements apply to the
full $W$ in \eqref{eq:normal-product}; the complementary
representation $U^\perp$ alone need not be a doubled quiver.
